%% annex/p03/maps-elementarity.tex --- Tier-2 proof compressions for v1.1
%% Part IV, source lines 6114--6485.  Every source-complete proof remains a
%% complete proof; the six items are ordered exactly as in the source chapter.

\begin{annexitem}{annex:p3.maps.1}{thm:p03-term-preservation}{Preservation of terms}
\begin{proof}
Structural induction on \(t\).  Variables are immediate, constants use
constant preservation, and \(f(t_1,\ldots,t_n)\) follows from the induction
hypotheses and then function preservation.
\end{proof}
\end{annexitem}

\begin{annexitem}{annex:p3.maps.2}{thm:p03-isomorphism-invariance}{Isomorphism invariance of satisfaction}
\begin{proof}
Induct on \(\varphi\).  Term preservation and injectivity handle equality;
term preservation and relation reflection handle relation atoms.  Boolean
connectives follow from the induction hypotheses.  For \(\exists y\,\psi\), a
witness \(a'\in M\) maps to \(h(a')\); conversely, surjectivity writes any
witness \(b\in N\) as \(h(a')\), and the induction hypothesis returns \(a'\).
The universal case uses the same witness argument on failure.
\end{proof}
\end{annexitem}

\begin{annexitem}{annex:p3.maps.3}{thm:p03-embedding-image}{Image of an embedding}
\begin{proof}
The homomorphism clauses give
\(c^{\mathcal N}=h(c^{\mathcal M})\in h(M)\) and
\(f^{\mathcal N}(h\bar a)=h(f^{\mathcal M}(\bar a))\in h(M)\), so the
restricted interpretations define \(h[\mathcal M]\subseteq\mathcal N\).
Indeed, the inclusion preserves and reflects atomic formulas because all its
interpretations are restrictions.
Corestricted to \(h(M)\), \(h\) is injective and surjective, preserves
constants and functions, and preserves and reflects each restricted relation.
It is therefore a surjective embedding, hence an isomorphism.
\end{proof}
\end{annexitem}

\begin{annexitem}{annex:p3.maps.4}{thm:p03-tarski-vaught}{Tarski--Vaught test}
\begin{proof}
If \(\mathcal M\preccurlyeq\mathcal N\), transfer the existential statement
to \(\mathcal M\), choose its witness there, and transfer its instance back.
Conversely assume the witness condition and induct on \(\psi(\bar x)\).
Atomic formulas agree because \(\mathcal M\subseteq\mathcal N\); Boolean cases
are immediate.  A witness in \(M\) transfers from \(\mathcal M\) to
\(\mathcal N\).  In the reverse existential direction the hypothesis chooses
\(b\in M\), and the induction hypothesis transfers
\(\varphi(b,\bar a)\) to \(\mathcal M\).  For a universal formula, instances
true in \(\mathcal N\) transfer to every \(b\in M\); if it fails in
\(\mathcal N\), apply the witness condition to
\(\exists y\,\neg\varphi(y,\bar a)\) and transfer the counterexample to
\(\mathcal M\).
\end{proof}
\end{annexitem}

\begin{annexitem}{annex:p3.maps.5}{thm:p03-diagram-tests}{Diagram tests for embeddings}
\begin{proof}
Given an embedding \(h\), interpret \(c_a\) in \(\mathcal N'\) as \(h(a)\);
preservation and reflection make its atomic diagram true.  Conversely define
\(h(a)=c_a^{\mathcal N'}\).  The sentences \(c_a\ne c_b\), the constant and
function equations, and the positive and negated relation instances prove
injectivity, preservation, and reflection.

For an elementary embedding the same expansion satisfies the elementary
diagram.  Conversely its atomic part yields an embedding.  For every
\(\varphi(\bar x)\) and \(\bar a\in M\), the elementary diagram contains
exactly the true member of
\(\{\varphi(\bar c_{\bar a}),\neg\varphi(\bar c_{\bar a})\}\); satisfaction
in \(\mathcal N'\) therefore gives the required biconditional, so \(h\) is
elementary.
\end{proof}
\end{annexitem}

\begin{annexitem}{annex:p3.maps.6}{thm:p03-elementary-chain}{Elementary chain theorem}
\begin{proof}
Induct on \(\psi(\bar x)\), simultaneously for every \(i<\delta\) and tuple
\(\bar a\in M_i\).  Atomic formulas agree by compatibility; Boolean cases
follow immediately.  A witness in \(M_i\) remains one in the union.  Conversely,
if \(b\in\bigcup_jM_j\) witnesses an existential in the union, choose
\(j\ge i\) with \(b\in M_j\); the induction hypothesis transfers its matrix
to \(\mathcal M_j\), and
\(\mathcal M_i\preccurlyeq\mathcal M_j\) transfers the existential back.
For a universal formula, transfer each union-true instance to
\(\mathcal M_i\); if it fails in the union, choose a counterexample in some
\(M_j\) with \(j\ge i\), transfer its matrix failure to \(\mathcal M_j\),
and then transfer the universal failure to \(\mathcal M_i\).
\end{proof}
\end{annexitem}
